\documentclass[12pt,a4paper]{ctexart}

% === 页面设置 ===
\usepackage[top=2.5cm, bottom=2.5cm, left=2.5cm, right=2.5cm]{geometry}
\usepackage{setspace}
\onehalfspacing

% === 常用宏包 ===
\usepackage{graphicx}
\usepackage{amsmath, amssymb}
\usepackage{booktabs}
\usepackage{tabularx}
\usepackage{multirow}
\usepackage{float}
\usepackage{caption}
\usepackage{subcaption}
\usepackage{listings}
\usepackage{xcolor}
\usepackage{hyperref}
\usepackage{fancyhdr}
\usepackage{titlesec}
\usepackage{enumitem}
\usepackage{tikz}
\usepackage{circuitikz}
\usepackage{pgfplots}
\pgfplotsset{compat=1.18}

% === 代码高亮设置 ===
\lstset{
    language=C,
    basicstyle=\small\ttfamily,
    keywordstyle=\color{blue}\bfseries,
    commentstyle=\color{green!60!black},
    stringstyle=\color{red},
    numbers=left,
    numberstyle=\tiny\color{gray},
    numbersep=8pt,
    frame=single,
    breaklines=true,
    breakatwhitespace=false,
    tabsize=4,
    showstringspaces=false,
    captionpos=b,
    aboveskip=10pt,
    belowskip=10pt
}

% === 页眉页脚 ===
\pagestyle{fancy}
\fancyhf{}
\fancyhead[L]{\small 全国大学生电子设计竞赛}
\fancyhead[R]{\small 设计报告}
\fancyfoot[C]{\thepage}
\renewcommand{\headrulewidth}{0.4pt}

% === 超链接设置 ===
\hypersetup{
    colorlinks=true,
    linkcolor=blue,
    citecolor=blue,
    urlcolor=blue
}

% === 图表编号格式 ===
\renewcommand{\thefigure}{\thesection-\arabic{figure}}
\renewcommand{\thetable}{\thesection-\arabic{table}}
\renewcommand{\theequation}{\thesection.\arabic{equation}}
\makeatletter
\@addreset{figure}{section}
\@addreset{table}{section}
\@addreset{equation}{section}
\makeatother

% ============================================================
\begin{document}

% === 封面 ===
\begin{titlepage}
\centering
\vspace*{2cm}

{\Huge\bfseries 全国大学生电子设计竞赛\\[0.3cm]设计报告}

\vspace{2cm}

{\LARGE\bfseries 【此处填写赛题名称】}

\vspace{1.5cm}

\begin{tabular}{rl}
{\large\bfseries 赛题编号：} & {\large \underline{\hspace{3cm}A题\hspace{3cm}}} \\[0.8cm]
{\large\bfseries 参赛学校：} & {\large \underline{\hspace{5cm}}} \\[0.8cm]
{\large\bfseries 队伍编号：} & {\large \underline{\hspace{5cm}}} \\[0.8cm]
{\large\bfseries 队员一：} & {\large \underline{\hspace{4cm}}（签名）} \\[0.8cm]
{\large\bfseries 队员二：} & {\large \underline{\hspace{4cm}}（签名）} \\[0.8cm]
{\large\bfseries 队员三：} & {\large \underline{\hspace{4cm}}（签名）} \\
\end{tabular}

\vfill
{\large 2025年\hspace{1cm}月\hspace{1cm}日}
\end{titlepage}

% === 目录 ===
\newpage
\tableofcontents
\newpage

% ============================================================
\section{系统方案论证}

\subsection{题目分析}

【在此分析赛题要求，列出基本要求和发挥部分，明确设计目标】

本题要求设计一个【简要描述系统功能】，主要技术指标如下：

\begin{table}[H]
\centering
\caption{主要技术指标}
\begin{tabular}{lll}
\toprule
\textbf{指标项} & \textbf{基本要求} & \textbf{发挥部分} \\
\midrule
指标1 & 值 & 值 \\
指标2 & 值 & 值 \\
指标3 & 值 & 值 \\
\bottomrule
\end{tabular}
\end{table}

\subsection{方案比较与选择}

\subsubsection{方案一：【方案名称】}

【描述方案原理、优缺点】

\subsubsection{方案二：【方案名称】}

【描述方案原理、优缺点】

\subsubsection{方案选择}

综合比较两种方案，选择方案【X】。理由如下：
\begin{enumerate}[label=(\arabic*)]
    \item 理由一
    \item 理由二
    \item 理由三
\end{enumerate}

% ============================================================
\section{理论分析与计算}

\subsection{【模块1理论分析】}

【在此进行必要的理论推导和参数计算】

关键公式如下：
\begin{equation}
    V_{out} = V_{in} \times \frac{R_2}{R_1 + R_2}
    \label{eq:voltage_divider}
\end{equation}

其中$V_{in}$为输入电压，$R_1$、$R_2$为分压电阻。

\subsection{【模块2理论分析】}

【理论推导内容】

\subsection{参数计算}

根据上述理论分析，计算关键参数：

\begin{table}[H]
\centering
\caption{关键参数计算结果}
\begin{tabular}{lll}
\toprule
\textbf{参数} & \textbf{计算值} & \textbf{实际取值} \\
\midrule
参数1 & 计算值 & 实际值 \\
参数2 & 计算值 & 实际值 \\
\bottomrule
\end{tabular}
\end{table}

% ============================================================
\section{电路与程序设计}

\subsection{系统总体框图}

\begin{figure}[H]
\centering
\begin{tikzpicture}[block/.style={draw, rectangle, minimum width=2.5cm, minimum height=1cm, align=center}]
    \node[block] (input) at (0,0) {输入\\调理};
    \node[block] (adc) at (3.5,0) {ADC\\采样};
    \node[block] (mcu) at (7,0) {MCU\\处理};
    \node[block] (dac) at (10.5,0) {DAC\\输出};
    \node[block] (output) at (14,0) {功率\\驱动};
    
    \draw[->] (input) -- (adc);
    \draw[->] (adc) -- (mcu);
    \draw[->] (mcu) -- (dac);
    \draw[->] (dac) -- (output);
\end{tikzpicture}
\caption{系统总体框图}
\label{fig:system_block}
\end{figure}

\subsection{硬件电路设计}

\subsubsection{【电路模块1】}

【电路设计说明，可插入原理图】

\begin{figure}[H]
\centering
% \includegraphics[width=0.8\textwidth]{schematic_module1.png}
\caption{【电路模块1】原理图}
\label{fig:circuit1}
\end{figure}

\subsubsection{【电路模块2】}

【电路设计说明】

\subsection{软件设计}

\subsubsection{程序总体流程}

\begin{figure}[H]
\centering
\begin{tikzpicture}[node distance=1.2cm,
    startstop/.style={draw, rounded rectangle, minimum width=2cm, align=center},
    process/.style={draw, rectangle, minimum width=2cm, minimum height=0.8cm, align=center},
    decision/.style={draw, diamond, minimum width=1.5cm, align=center}]
    
    \node[startstop] (init) {系统初始化};
    \node[process, below=of init] (read) {读取传感器};
    \node[process, below=of read] (process) {数据处理};
    \node[process, below=of process] (output) {输出控制};
    \node[decision, below=of output] (check) {是否停止？};
    
    \draw[->] (init) -- (read);
    \draw[->] (read) -- (process);
    \draw[->] (process) -- (output);
    \draw[->] (output) -- (check);
    \draw[->] (check) -| node[anchor=south] {否} ++(3,0) |- (read);
    \draw[->] (check) -- node[anchor=east] {是} ++(0,-1) node[below] {结束};
\end{tikzpicture}
\caption{主程序流程图}
\label{fig:flowchart}
\end{figure}

\subsubsection{关键代码}

\begin{lstlisting}[caption={关键算法实现}, label={lst:key_code}]
/* 关键代码示例 */
void System_Process(void) {
    // 读取ADC
    uint16_t adc_val = ADC_Read();
    
    // 数据处理
    float result = Algorithm_Process(adc_val);
    
    // 输出
    DAC_Output(result);
}
\end{lstlisting}

% ============================================================
\section{测试方案与测试结果}

\subsection{测试仪器}

\begin{table}[H]
\centering
\caption{测试仪器列表}
\begin{tabular}{lll}
\toprule
\textbf{仪器名称} & \textbf{型号} & \textbf{主要参数} \\
\midrule
数字示波器 & Tektronix TDS2022 & 200MHz, 2GS/s \\
信号发生器 & AFG3022 & 25MHz, 双通道 \\
数字万用表 & Fluke 17B & 4000字 \\
直流稳压电源 & \text{-} & 0\textasciitilde 30V, 5A \\
\bottomrule
\end{tabular}
\end{table}

\subsection{测试方案}

【描述测试方法和步骤】

\subsection{测试数据}

\begin{table}[H]
\centering
\caption{【测试项目1】测试数据}
\begin{tabular}{cccc}
\toprule
\textbf{输入} & \textbf{理论值} & \textbf{实测值} & \textbf{误差} \\
\midrule
1 & 1.00 & 0.99 & 1.0\% \\
2 & 2.00 & 2.01 & 0.5\% \\
3 & 3.00 & 2.98 & 0.7\% \\
\bottomrule
\end{tabular}
\end{table}

\subsection{测试结果分析}

\begin{itemize}
    \item 基本要求完成情况：【逐条说明】
    \item 发挥部分完成情况：【逐条说明】
    \item 测试结论：【总结性评价】
\end{itemize}

% ============================================================
\section{总结}

\subsection{系统特点}

\begin{enumerate}
    \item 特点一
    \item 特点二
    \item 特点三
\end{enumerate}

\subsection{不足与改进}

\begin{enumerate}
    \item 不足一及改进方向
    \item 不足二及改进方向
\end{enumerate}

% ============================================================
\section*{参考文献}
\begin{enumerate}[label={[\arabic*]}]
    \item 参考文献1
    \item 参考文献2
    \item 参考文献3
\end{enumerate}

% === 附录 ===
\appendix
\section{完整电路原理图}

% \includegraphics[width=\textwidth]{full_schematic.png}

\section{完整程序代码}

\begin{lstlisting}[language=C, caption={main.c}]
/* 完整主程序代码 */
#include "stm32f10x.h"

int main(void) {
    SystemInit();
    // ... 完整代码
    while(1) {
        // 主循环
    }
}
\end{lstlisting}

\end{document}
